\section{Tensors and arithmetic exponents}\label{sec:preliminaries}

We first fix the tensor conventions shared by the square and rectangular
constructions, then relate rank and explicit polynomial degenerations to
arithmetic exponents.  The equal-summand inequality at the end of the
section will convert separated tensors into rectangular bounds.

All vector spaces and tensors are over \(\C\).  A tensor on three finite
coordinate sets \(X,Y,Z\) is identified with a trilinear form
\[
 T=\sum_{x,y,z}T_{xyz}\,x y z.
\]
A restriction is obtained by a linear map on each of the three sides.
In particular, a local zeroing-out sets selected variables on each side
to zero.  A degeneration \(U\unrhd V\) permits these maps to depend
polynomially or Laurent polynomially on a parameter and extracts the
transformed tensor's first nonzero coefficient, after an overall rescaling.
The maps may be composed finitely many times.
Finitely many parameters can be replaced by sufficiently separated powers
of a single parameter.

The rank \(\rank(T)\) is the minimum number of products of three linear
forms whose sum is \(T\).  Its border rank \(\brank(T)\) is the least
integer \(r\) such that tensors of rank at most \(r\) can approach
\(T\).  Border rank is submultiplicative under tensor product and cannot
increase under degeneration.  The direct sum \(U\oplus V\) uses
disjoint variable spaces on all three sides.  A claimed direct sum must
therefore have both disjoint pieces and no additional terms joining their
variables.

Our matrix-multiplication convention is
\[
 \langle a,b,c\rangle
 =\sum_{i=1}^{a}\sum_{j=1}^{b}\sum_{k=1}^{c}x_{ij}y_{jk}z_{ki},
 \qquad \Vol\langle a,b,c\rangle=abc.
\]
For positive integers \(l_1,l_2,l_3\), write
\(\mathcal R(l_1,l_2,l_3)=\rank\langle l_1,l_2,l_3\rangle\).

For \(a_1,a_2,a_3\ge0\), let \(W(a_1,a_2,a_3)\) denote the arithmetic
exponent for the three dimensions \(\lceil n^{a_i}\rceil\).
The trilinear interpretation makes \(W\) invariant under permutations,
including those implemented by transposition.
In particular \(w(k)=W(1,k,1)\).

\begin{lemma}\label{lem:rank-exponent}
For \(a_i\ge0\),
\begin{equation}\label{eq:rank-exponent}
 W(a_1,a_2,a_3)=
 \lim_{s\to\infty}\frac{\log\mathcal R
   (\lceil e^{sa_1}\rceil,\lceil e^{sa_2}\rceil,
          \lceil e^{sa_3}\rceil)}{s}.
\end{equation}
Moreover, \(W\) is coordinatewise nondecreasing, homogeneous,
subadditive, permutation invariant, and continuous.  It satisfies
\[
 W(a_1,a_2,a_3)\ge\max_{i\ne j}(a_i+a_j).
\]
\end{lemma}

\begin{proof}
Fix \(s>0\), put \(l_i=\lceil e^{sa_i}\rceil\), and take a rank-\(r\)
scheme for this triple.  Each of the three flattenings has rank
\(l_i l_j\), so \(r\ge\max_{i\ne j}l_i l_j\).
The scheme applied recursively to blocks of dimensions \(l_i^{h-1}\)
uses \(r\) smaller products and at most a fixed constant times
\((\max_{i\ne j}l_i l_j)^{h-1}\) linear operations at each level.
Its total cost is \(O((h+1)r^h)\).
All products keep a left-input block before a right-input block, so
rectangular block multiplication is legitimate.
Padding with \(h=\lceil t/s\rceil\) proves the corresponding arithmetic
upper bound for arbitrary large dimensions.
The same padding and tensor-power argument shows that the limsup in
\eqref{eq:rank-exponent} is at most the infimum over fixed \(s>0\).
Hence the rank limit exists.

Conversely, a straight-line arithmetic computation of the bilinear
matrix product gives rank at most a constant times its operation count.
Introduce separate formal variables for the two operands and retain
bidegrees at most \((1,1)\).
Each gate introduces only a bounded number of new products of a linear
form in the first operand and one in the second; the remaining mixed
coefficients are linear combinations of earlier ones.
For divisions we use Strassen's division-elimination method
\cite{strassen-division}, in the following bilinear adaptation.
If divisions are allowed, first translate to a generic constant input
at which every denominator is nonzero and use the same truncated-series
argument.  The mixed coefficient of the output is still the desired
bilinear map.
This proves the converse inequality in \eqref{eq:rank-exponent}.

Restriction and padding give monotonicity; tensor products give
subadditivity; rescaling \(s\) gives homogeneity.
These statements include coordinates equal to zero.
For continuity, compare two triples through their coordinatewise
maximum and use the naive bound on the nonnegative increment:
the change in \(W\) is at most the sum of absolute coordinate changes.
Finally the flattening bounds give the displayed lower bound.
\end{proof}

All logarithms in this paper are natural.
The properties in \cref{lem:rank-exponent} imply that \(w\) is convex.
For instance, subadditivity and homogeneity applied to
\(\lambda(1,k_1,1)+(1-\lambda)(1,k_2,1)\) give its usual convexity
inequality.

The approximate bilinear algorithms of \citet{bini} and the
asymptotic-sum method of \citet{schonhage} motivate the following
explicit form of a rank budget.  We use polynomial border-rank witnesses.
Saying that \(T\) has \emph{cost} \(R\) means that \(T\), up to a nonzero
scalar, is the first nonzero coefficient of a polynomial or Laurent
polynomial tensor family with a rank-\(R\) decomposition.
This is a supplied witness and a certified upper bound, which need not
be minimal.  It implies \(\brank(T)\le R\); it is not our definition of
border rank.  All uses of cost below have an explicit witness, so no
converse assertion about arbitrary border-rank limits is needed.
Coordinate restrictions and monomial leading-term degenerations preserve
such witnesses and their budgets.  Finitely many parameters can be
combined into one by assigning sufficiently separated powers, so a
composition of the operations used below again has a univariate witness.

The following interpolation argument is the border-to-exact transfer
for polynomial approximate bilinear algorithms \cite{bini}.

\begin{lemma}\label{lem:border-overhead}
For a fixed tensor of cost \(R\), its \(N\)-th tensor power has exact
rank at most \(O(N+1)R^N\), where the implicit constant may depend on
the fixed witness.
\end{lemma}

\begin{proof}
Clear negative parameter powers.
The degree of the \(N\)-th power of the witness is \(O(N)\), and its
desired coefficient is the \(N\)-th tensor power of the first
coefficient.  Univariate interpolation expresses that coefficient using
\(O(N+1)\) evaluations, each of rank at most \(R^N\).
\end{proof}

A dot tensor on the pair \(Y,Z\) of size \(D\) is
\(x_0\sum_{j=1}^D y_jz_j\).
The product of dot tensors on the three pairs is a matrix-multiplication
tensor.  More precisely, if their lengths on \(XY,XZ,YZ\) are
\(d_{XY},d_{XZ},d_{YZ}\), respectively, their product is
\[
 \langle d_{XZ},d_{XY},d_{YZ}\rangle.
\]
Thus the edge \(XZ\) supplies the index \(i\), the edge \(XY\)
supplies \(j\), and the edge \(YZ\) supplies \(k\) in our convention.
Dot sizes on the same pair multiply.  A tuple of logarithmic edge rates
listed in the order \((XY,XZ,YZ)\) therefore gives matrix-dimension
rates in the order \((XZ,XY,YZ)\).  By the permutation invariance in
Lemma~\ref{lem:rank-exponent}, either order gives the same value of
\(W\).  This accounts for the edge-rate order used in the later
conditional-label entropy bound.

\begin{lemma}[Equal rectangular summands]\label{lem:rectangular-sum}
Suppose a tensor of cost \(R\) restricts or degenerates to a direct sum
of \(M\) copies of the matrix-multiplication tensor with dimensions
\(l_1,l_2,l_3\).  Then
\begin{equation}\label{eq:rectangular-sum}
 W(\log l_1,\log l_2,\log l_3)\le \log R-\log M.
\end{equation}
\end{lemma}

\begin{proof}
Write the left side as \(u\).
If \(u=0\), a flattening of the \(N\)-th power of the direct sum,
together with \cref{lem:border-overhead}, gives
\(M^N\le O(N+1)R^N\); thus \(R\ge M\).

Otherwise choose an integer \(c\) with \(cu>\log M\).
A rank scheme for dimensions \(l_i^{cN}\) can be interpreted as requests
for products of blocks with dimensions \(l_i^N\).
The \(N\)-th power of the direct sum supplies \(M^N\) such requests in
one batch, at cost \(O(N+1)R^N\).
The requests may depend on common inputs: substituting those linear
forms into the independent batch inputs is a restriction.
Pad the last batch if necessary.  Consequently
\[
 \mathcal R(l_1^{(1+c)N},l_2^{(1+c)N},l_3^{(1+c)N})
 \le O(N+1)R^N
 \left\lceil
 \frac{\mathcal R(l_1^{cN},l_2^{cN},l_3^{cN})}{M^N}
 \right\rceil .
\]
By \cref{lem:rank-exponent}, the expression inside the ceiling grows
exponentially because \(cu>\log M\).
Take logarithms, divide by \(N\), and pass to the limit to obtain
\((1+c)u\le \log R+cu-\log M\), proving the claim.
\end{proof}

This is the equal-summand rectangular form of the asymptotic-sum
principle \cite{schonhage}.
The proof applies it to each fixed instance before varying that
instance.  Later auxiliary tensors may therefore grow with a type
parameter without any uniformity assumption on the interpolation
constant in \cref{lem:border-overhead}.
